% -*- coding: utf-8 -*-
% !TEX program = xelatex
\documentclass[plain,noanswer,medmath]{../../template/jnuexam}
\usepackage{../../template/kysx}

\begin{document}


\examtitle{2006}{1}


\makepart{填空题}{1～6小题，每小题4分，共24分}


\begin{problem}
$\lim_{x \to 0} \frac{x\ln(1 + x)}{1 - \cos x} =$ \fillin{}．
\end{problem}


\begin{problem}
微分方程$y' = \frac{y(1 - x)}{x}$的通解是 \fillin{}．
\end{problem}


\begin{problem}
设$\Sigma$是锥面$z = \sqrt{x^2 + y^2}$（$0 \le z \le 1$）的下侧，则\goodbreak
$\iint_{\Sigma}x\dy\dz + 2y\dz\dx + 3(z - 1)\dx\dy =$ \fillin{}．
\end{problem}


\begin{problem}
点$(2,1,0)$到平面$3x + 4y + 5z = 0$的距离$d=$ \fillin{}．
\end{problem}


\begin{problem}
设矩阵$A = \begin{pmatrix}
   2 & 1 \\
  -1 & 2
\end{pmatrix}$，$E$为$2$阶单位矩阵，矩阵$B$满足$BA = B + 2E$，则$|B| =$ \fillin{}．
\end{problem}


\begin{problem}
设随机变量$X$与$Y$相互独立，且均服从区间$[0,3]$上的均匀分布，则\goodbreak
$P\left\{\max \{X,\;Y\} \le 1 \right\} =$ \fillin{}．
\end{problem}


\makepart{选择题}{7～14小题，每小题4分，共32分}


\begin{problem}
设函数$y = f(x)$具有二阶导数，且$f'(x) > 0,\;f''(x) > 0$，
$\Delta x$为自变量$x$在$x_0$处的增量，$\Delta y$与$\dy$分别为$f(x)$在点$x_0$处对应的增量与微分，
若$\Delta x > 0$，则\pickout{}
\begin{abcd}
\item $0 < \dy < \Delta y$．
\item $0 < \Delta y < \dy$．
\item $\Delta y < \dy < 0$．
\item $\dy < \Delta y < 0$．
\end{abcd}
\end{problem}


\begin{problem}
设$f(x,y)$为连续函数，则$\int_0^{\frac{\pi}{4}}\d\theta\int_0^1 f(r\cos\theta,r\sin\theta)r\dr$等于\pickout{}
\begin{abcd}
\item $\int_0^{\frac{\sqrt{2}}{2}} \dx \int_x^{\sqrt{1 - x^2}} f(x,y)\dy$．
\item $\int_0^{\frac{\sqrt{2}}{2}} \dx \int_0^{\sqrt{1 - x^2}} f(x,y)\dy$．
\item $\int_0^{\frac{\sqrt{2}}{2}} \dy \int_y^{\sqrt{1 - y^2}} f(x,y)\dx$．
\item $\int_0^{\frac{\sqrt{2}}{2}} \dy \int_0^{\sqrt{1 - y^2}} f(x,y)\dx$．
\end{abcd}
\end{problem}


\begin{problem}
若级数$\sum_{n = 1}^{\infty}a_n$收敛，则级数\pickout{}
\begin{abcd}
\item $\sum_{n = 1}^{\infty}\left| a_n \right|$收敛．
\item $\sum_{n = 1}^{\infty}(- 1)^n a_n$收敛．
\item $\sum_{n = 1}^{\infty}a_n a_{n + 1}$收敛．
\item $\sum_{n = 1}^{\infty}\frac{a_n + a_{n + 1}}{2}$收敛．
\end{abcd}
\end{problem}


\begin{problem}
设$f(x,y)$与$\varphi(x,y)$均为可微函数，且$\varphi'_y(x,y) \ne 0$．
已知$(x_0,y_0)$是$f(x,y)$在约束条件$\varphi (x,y) = 0$下的一个极值点，
下列选项正确的是\pickout{}
\begin{abcd}
\item 若$f'_x(x_0,y_0) = 0$，则$f'_y(x_0,y_0) = 0$．
\item 若$f'_x(x_0,y_0) = 0$，则$f'_y(x_0,y_0) \ne 0$．
\item 若$f'_x(x_0,y_0) \ne 0$，则$f'_y(x_0,y_0) = 0$．
\item 若$f'_x(x_0,y_0) \ne 0$，则$f'_y(x_0,y_0) \ne 0$．
\end{abcd}
\end{problem}


\begin{problem}
设$a_1,\;a_2,\; \cdots ,\;a_s$均为$n$维列向量，$A$是$m \times n$矩阵，下列选项正确的是\pickout{}
\begin{abcd}
\item 若$a_1,\;a_2,\; \cdots ,\;a_s$线性相关，则$Aa_1,\;Aa_2, \cdots ,\;Aa_s$线性相关．
\item 若$a_1,\;a_2,\; \cdots ,\;a_s$线性相关，则$Aa_1,\;Aa_2, \cdots ,\;Aa_s$线性无关．
\item 若$a_1,\;a_2,\; \cdots ,\;a_s$线性无关，则$Aa_1,\;Aa_2, \cdots ,\;Aa_s$线性相关．
\item 若$a_1,\;a_2,\; \cdots ,\;a_s$线性无关，则$Aa_1,\;Aa_2, \cdots ,\;Aa_s$线性无关．
\end{abcd}
\end{problem}


\begin{problem}
设$A$为$3$阶矩阵，将$A$的第$2$行加到第$1$行得$B$，
再将$B$的第$1$列的$- 1$倍加到第$2$列得$C$，
记$P = \begin{pmatrix}
  1 & 1 & 0 \\
  0 & 1 & 0 \\
  0 & 0 & 1
\end{pmatrix}$，则\pickout{}
\begin{abcd}
\item $C = P^{-1} AP$．
\item $C = PA P^{-1}$．
\item $C = P^T AP$．
\item $C = PA P^T$．
\end{abcd}
\end{problem}


\begin{problem}
设$A,\;B$为随机事件，且$P(B) > 0,\;P(A|B) = 1$，则必有\pickout{}
\begin{abcd}
\item $P(A \cup B) > P(A)$．
\item $P(A \cup B) > P(B)$．
\item $P(A \cup B) = P(A)$．
\item $P(A \cup B) = P(B)$．
\end{abcd}
\end{problem}


\begin{problem}
设随机变量$X$服从正态分布$N(\mu_1,\;\sigma_1^2)$，$Y$服从正态分布$N(\mu_2,\;\sigma_2^2)$，
且\goodbreak$P\{|X - \mu_1| < 1\} > P\{|Y - \mu_2| < 1\}$，则必有\pickout{}
\begin{abcd}
\item $\sigma_1 < \sigma_2$．
\item $\sigma_1 > \sigma_2$．
\item $\mu_1 < \mu_2$．
\item $\mu_1 > \mu_2$．
\end{abcd}
\end{problem}


\makepart{解答题}{15～23小题，共94分}


\begin{problem}[points=10]
设区域$D = \left\{\left(x,y \right)\left| x^2 + y^2 \le 1,x \ge 0 \right. \right\}$，
计算二重积分$I = \iint_D \frac{1 + xy}{1 + x^2 + y^2}\dx\dy$．
\end{problem}


\begin{problem}[points=12]
设数列$\{x_n\}$满足$0<x_1<\pi$, $x_{n+1} = \sin x_n\;(n = 1,2, \cdots)$．\par
\begin{enumerate*}
\item 证明$\lim_{n \to \infty} x_n$存在，并求该极限；
\item 计算$\lim_{n \to \infty} \left(\frac{x_{n + 1}}{x_n}\right)^{\frac{1}{x_n^2}}$．
\end{enumerate*}
\end{problem}


\begin{problem}[points=12]
将函数$f(x) = \frac{x}{2 + x - x^2}$展开成$x$的幂级数．
\end{problem}


\begin{problem}[points=12]
设函数$f(u)$在$(0,+\infty)$内具有二阶导数，
且$z = f\left(\sqrt{x^2 + y^2}\right)$满足等式
$$\frac{\pd^2z}{\pdx^2} + \frac{\pd^2z}{\pdy^2}=0.$$
\begin{enumerate*}
\item 验证$f''(u) + \frac{f'(u)}{u} = 0$．
\item 若$f(1) = 0,f'(1) = 1$，求函数$f(u)$的表达式．
\end{enumerate*}
\end{problem}


\begin{problem}[points=12]
设在上半平面$D = \left\{(x,y) \big| y > 0\right\}$内，函数$f(x,y)$是有连续偏导数，
且对任意的$t > 0$都有$f(tx,ty) = t^{-2} f(x,y)$．
证明：对$D$内的任意分段光滑的有向简单闭曲线$L$，
都有$\oint_L y f(x,y)\dx - xf(x,y)\dy = 0$．
\end{problem}


\begin{problem}[points=9]
已知非齐次线性方程组$\begin{cases}
  x_1 + x_2 + x_3 + x_4 = - 1 \\
  4x_1 + 3x_2 + 5x_3 - x_4 = - 1 \\
  ax_1 + x_2 + 3x_3 + bx_4 = 1
\end{cases}$有$3$个线性无关的解．
\begin{enumerate}
\item 证明方程组系数矩阵$A$的秩$r(A) = 2$；
\item 求$a,b$的值及方程组的通解．
\end{enumerate}
\end{problem}


\begin{problem}[points=9]
设$3$阶实对称矩阵$A$的各行元素之和均为$3$，
向量$\alpha_1 = \left(- 1,2, - 1 \right)^T$,$\alpha_2 = \left(0, - 1,1 \right)^T$
是线性方程组$Ax = 0$的两个解．
\begin{enumerate}
\item 求$A$的特征值与特征向量；
\item 求正交矩阵$Q$和对角矩阵$\Lambda$，使得$Q^T AQ = \Lambda$．
\end{enumerate}
\end{problem}


\begin{problem}[points=9]
随机变量$X$的概率密度为
$$f_X(x) = \begin{cases}
  \frac{1}{2}, & -1 < x < 0; \\
  \frac{1}{4}, & 0 \le x < 2; \\
            0, & \text{其他}.
\end{cases}$$
$\text{令}Y = X^2$，$F\left(x,y \right)$为二维随机变量$(X,Y)$的分布函数．求\par
\begin{enumerate*}
\item $Y$的概率密度$f_Y(y)$；
\item $F\left(-\frac{1}{2},4\right)$．
\end{enumerate*}
\end{problem}


\begin{problem}[points=9]
设总体$X$的概率密度为
$$f(x,\theta) = \begin{cases}
      \theta , & 0 < x < 1, \\
  1 - \theta , & 1 \le x < 2, \\
            0, & \text{其它},
\end{cases}$$
其中$\theta$是未知参数（$0 < \theta < 1$），
$X_1,X_2\cdots ,X_{n}$为来自总体$X$的简单随机样本，
记$N$为样本值$x_1,x_2\cdots ,x_n$中小于$1$的个数，求$\theta$的最大似然估计．
\end{problem}


\end{document}
